% 第七章
\chapter{磁群与共表示}

% 7.1 Shubnikov 群及其推导
\section{Shubnikov 群及其推导}

230 个空间群（Fedorov 群）已在第一章讨论过，并在 3.5 节给出了详细的分类。在 Shubnikov 于 1951 年引入反对称操作的概念之前，这些群一直被视为晶体对称性研究的最高成就。通过在晶体的通常位置坐标之外引入一个只有两个可能取值的额外坐标——例如颜色（黑或白）、符号（$+$ 或 $-$）、或磁矩方向（平行或反平行于给定方向）——可以定义\textit{反对称操作}（operation of anti-symmetry）：即使该坐标改变值的操作，例如把黑变为白、白变为黑。若反对称操作记作 $\mathcal{R}$，则可以有复合的反对称操作，它对应于同时施行一个通常的点群或空间群操作 $\{R \mid \mathbf{v}\}$ 与反对称操作 $\mathcal{R}$；这类似于由纯转动操作与反演操作构成的非本征点群操作。允许这类操作后，便可以导出黑白群或\textit{磁点群}（magnetic point groups），共 58 个；以及黑白群或\textit{磁空间群}（magnetic space groups），共 1191 个。磁点群由 Shubnikov (1951) 导出，并由 Tavger and Zaitsev (1956) 方便地列表；磁空间群由 Zamorzaev 在 1953 年的学位论文以及 Belov, Neronova, and Smirnova (1955) 导出。二维黑白空间群曾由 Cochran (1952) 讨论。如果再把“全白群”与“灰群”包括进来，反对称操作的引入导致 $58 + 32 + 32 = 122$ 个点群与 $1191 + 230 + 230 = 1651$ 个空间群。本章使用术语 \textit{Shubnikov 群}统称所有这些群，从上下文可以清楚所指的是点群还是空间群。这些群的大量重要早期俄文文献现已译成英文（Shubnikov and Belov 1964），很容易获得。Birss (1963, 1964) 对晶体各种张量形式与磁点群的关系给出了详尽的叙述。

% 7.2 Shubnikov 群的分类
\section{Shubnikov 群的分类}

先考虑点群。若 $\mathbf{G}$ 是本书前面各节讨论过的普通点群之一，则有三种类型的 Shubnikov 群与之对应，记为 I、II 或 III 型：

I 型：普通点群（32 个）；

II 型：灰点群（grey point groups，32 个）；

III 型：黑白群或磁点群（58 个）。

括号中的数字给出每种类型点群的数目，这些 Shubnikov 点群总数为 122。

I 型群中不存在反对称操作 $\mathcal{R}$；这些群 $\mathbf{G}$ 已在表 1.3 中列出，这里不再多说。

我们引入的、允许取两个值之一的额外坐标 $s$，在 II 型群中被假定同时取两个值，于是 $\mathbf{G}$ 的任何操作使 $s$ 不变，而 $\mathcal{R}$ 乘以 $\mathbf{G}$ 的任何操作则把黑变白、白变黑，从而也使 $s$ 不变。于是 $\mathcal{R}$ 乘以 $\mathbf{G}$ 的任何操作也是该群的元素。

III 型群（黑白群）含有 $\mathcal{R}$ 与一半 $\mathbf{G}$ 的操作之积。这三类点群的完整清单连同其生成元在原书表 7.1 中给出；表 7.2 与 7.3 列出 Shubnikov 空间群的分类与生成元；图 7.1--7.4 示意方块的对称操作、黑白方格、黑白布拉维格子等。以下为原书第 590--617 页的扫描。

\begin{figure}[H]
\centering
\includegraphics[width=0.86\textwidth]{c7_p603.jpg}
\caption*{\textbf{表 7.1}--7.4 与图 7.1--7.4：Shubnikov 群的推导与分类（原书第 590--617 页扫描）}
\end{figure}
\begin{figure}[H]
\centering
\includegraphics[width=0.86\textwidth]{c7_p604.jpg}
\caption*{原书第 591 页}
\end{figure}
\begin{figure}[H]
\centering
\includegraphics[width=0.86\textwidth]{c7_p605.jpg}
\caption*{原书第 592 页}
\end{figure}
\begin{figure}[H]
\centering
\includegraphics[width=0.86\textwidth]{c7_p606.jpg}
\caption*{原书第 593 页}
\end{figure}
\begin{figure}[H]
\centering
\includegraphics[width=0.86\textwidth]{c7_p607.jpg}
\caption*{原书第 594 页}
\end{figure}
\begin{figure}[H]
\centering
\includegraphics[width=0.86\textwidth]{c7_p608.jpg}
\caption*{原书第 595 页}
\end{figure}
\begin{figure}[H]
\centering
\includegraphics[width=0.86\textwidth]{c7_p609.jpg}
\caption*{原书第 596 页}
\end{figure}
\begin{figure}[H]
\centering
\includegraphics[width=0.86\textwidth]{c7_p610.jpg}
\caption*{原书第 597 页}
\end{figure}
\begin{figure}[H]
\centering
\includegraphics[width=0.86\textwidth]{c7_p611.jpg}
\caption*{原书第 598 页}
\end{figure}
\begin{figure}[H]
\centering
\includegraphics[width=0.86\textwidth]{c7_p612.jpg}
\caption*{原书第 599 页}
\end{figure}
\begin{figure}[H]
\centering
\includegraphics[width=0.86\textwidth]{c7_p613.jpg}
\caption*{原书第 600 页}
\end{figure}
\begin{figure}[H]
\centering
\includegraphics[width=0.86\textwidth]{c7_p614.jpg}
\caption*{原书第 601 页}
\end{figure}
\begin{figure}[H]
\centering
\includegraphics[width=0.86\textwidth]{c7_p615.jpg}
\caption*{原书第 602 页}
\end{figure}
\begin{figure}[H]
\centering
\includegraphics[width=0.86\textwidth]{c7_p616.jpg}
\caption*{原书第 603 页}
\end{figure}
\begin{figure}[H]
\centering
\includegraphics[width=0.86\textwidth]{c7_p617.jpg}
\caption*{原书第 604 页}
\end{figure}
\begin{figure}[H]
\centering
\includegraphics[width=0.86\textwidth]{c7_p618.jpg}
\caption*{原书第 605 页}
\end{figure}
\begin{figure}[H]
\centering
\includegraphics[width=0.86\textwidth]{c7_p619.jpg}
\caption*{原书第 606 页}
\end{figure}
\begin{figure}[H]
\centering
\includegraphics[width=0.86\textwidth]{c7_p620.jpg}
\caption*{原书第 607 页}
\end{figure}
\begin{figure}[H]
\centering
\includegraphics[width=0.86\textwidth]{c7_p621.jpg}
\caption*{原书第 608 页}
\end{figure}
\begin{figure}[H]
\centering
\includegraphics[width=0.86\textwidth]{c7_p622.jpg}
\caption*{原书第 609 页}
\end{figure}
\begin{figure}[H]
\centering
\includegraphics[width=0.86\textwidth]{c7_p623.jpg}
\caption*{原书第 610 页}
\end{figure}
\begin{figure}[H]
\centering
\includegraphics[width=0.86\textwidth]{c7_p624.jpg}
\caption*{原书第 611 页}
\end{figure}
\begin{figure}[H]
\centering
\includegraphics[width=0.86\textwidth]{c7_p625.jpg}
\caption*{原书第 612 页}
\end{figure}
\begin{figure}[H]
\centering
\includegraphics[width=0.86\textwidth]{c7_p626.jpg}
\caption*{原书第 613 页}
\end{figure}
\begin{figure}[H]
\centering
\includegraphics[width=0.86\textwidth]{c7_p627.jpg}
\caption*{原书第 614 页}
\end{figure}
\begin{figure}[H]
\centering
\includegraphics[width=0.86\textwidth]{c7_p628.jpg}
\caption*{原书第 615 页}
\end{figure}
\begin{figure}[H]
\centering
\includegraphics[width=0.86\textwidth]{c7_p629.jpg}
\caption*{原书第 616 页}
\end{figure}
\begin{figure}[H]
\centering
\includegraphics[width=0.86\textwidth]{c7_p630.jpg}
\caption*{原书第 617 页}
\end{figure}

% 7.3 反幺正操作与磁群的共表示
\section{反幺正操作与磁群的共表示}

属于 I 型的 Shubnikov 群的表示已在第三、四章讨论过，并在第五、六章详细给出。通过引入反对称操作 $\mathcal{R}$，我们导出了一整套新群：II、III、IV 型 Shubnikov 群。由式 (7.2.1) 与 (7.2.3) 定义的 II 型 Shubnikov 群是直积群。如果把 $\mathcal{R}$（被视为反转磁矩方向的操作）当作幺正算符，确定这些群的表示并没有太大困难。但在研究磁对称性时，算符 $\mathcal{R}$ 是时间反演操作，记作 $\theta$，它是\textit{反幺正}的（anti-unitary）。因此必须能把表示理论推广到包含反幺正算符的群。该理论最初由 Wigner 于 1932 年给出，最容易读到的表述是 Wigner 经典著作的英译本（Wigner 1959，第 26 章）。理论的发展还由 Dimmock and Wheeler (1962\textit{a}, 1962\textit{b})（另见 Dimmock (1963\textit{e})）、Bradley and Davies (1968) 以及 Chaldyshev, Kudryavtseva, and Karavaev 的一系列论文讨论过。幺正群有表示与不可约表示；而对非幺正群——一半元素为幺正、另一半为反幺正的群——则有\textit{共表示}（corepresentations）与不可约共表示。本节陈述共表示的基本理论。对不想追完全部详细证明的读者，本节末尾（从式 (7.3.45) 起）总结了主要结果。

设 $\mathbf{M}$ 是一个非幺正群，$\mathbf{G}$ 为其幺正子群（指数为 2），$\mathbf{M} = \mathbf{G} + (\theta R_{0})\mathbf{G}$。设 $\Gamma^{i}$ 是 $\mathbf{G}$ 的不可约表示。Wigner 证明了：$\mathbf{M}$ 的不可约共表示由 $\Gamma^{i}$ 按\textit{三种情形}导出：

\noindent（a）若 $\Gamma^{i}$ 与 $\Gamma^{i*}(A^{-1}RA)$ 不等价（“普通”情形），则 $D\Gamma^{i}$ 由 $\Gamma^{i} \uparrow \mathbf{M}$ 给出，维数加倍；

\noindent（b）若 $\Gamma^{i} \simeq \Gamma^{i*}(A^{-1}RA)$ 但 $\Gamma^{i}$ 可取成实（或可由等价变换化为实），则 $D\Gamma^{i}$ 由 $\Gamma^{i}$ 的对称化平方构造；

\noindent（c）若 $\Gamma^{i} \simeq \Gamma^{i*}(A^{-1}RA)$ 而 $\Gamma^{i}$ 只能取成复的，则 $D\Gamma^{i}$ 由 $\Gamma^{i}$ 加上其“伴随”构造。

判别属于 (a)、(b) 还是 (c) 的检验由式 (7.3.x) 给出，即计算 $\sum_{R}\chi^{i}(R^{2})$（情形 (b)/(c) 的判据）；这与 4.6 节关于时间反演的实性理论完全平行。本节还给出共表示的特征标与矩阵代表的构造方法，并从式 (7.3.45) 起总结了全部主要结果。以下为表 7.5--7.7 的扫描。

\begin{figure}[H]
\centering
\includegraphics[width=0.86\textwidth]{c7_p641.jpg}
\caption*{\textbf{表 7.5}--7.7：共表示的基本表示与因子系统（原书第 628--630 页扫描）}
\end{figure}
\begin{figure}[H]
\centering
\includegraphics[width=0.86\textwidth]{c7_p642.jpg}
\caption*{原书第 629 页}
\end{figure}
\begin{figure}[H]
\centering
\includegraphics[width=0.86\textwidth]{c7_p643.jpg}
\caption*{原书第 630 页}
\end{figure}

% 7.4 共表示的 Kronecker 积
\section{共表示的 Kronecker 积}

本节考虑非幺正群 $\mathbf{M}$ 的任意两个不可约共表示的内 Kronecker 积的定义与约化。沿用 7.3 节的记号，以 $\Gamma^{i}$ 记幺正子群 $\mathbf{G}$ 的不可约表示，并为方便起见定义
\begin{equation*}
\Gamma^{i'}(R) = \overline{\Gamma}^{i}(R) = \Gamma^{i*}(A^{-1}RA).
\tag{7.4.1}
\end{equation*}
以 $D\Gamma^{i}$ 记由 $\Gamma^{i}$ 导出的 $\mathbf{M}$ 的不可约共表示。

我们假定幺正子群内部的内 Kronecker 积的分解已知，即式
\begin{equation*}
\Gamma^{i} \boxtimes \Gamma^{j} \equiv \sum_{k}c_{ij, k}\Gamma^{k}
\tag{7.4.2}
\end{equation*}
中的 Clebsch--Gordan 系数 $c_{ij, k}$ 已知。记 $\psi^{i}$ 为 $\Gamma^{i}$ 的特征标，由式 (1.3.24)，
\begin{equation*}
c_{ij, k} = \frac{1}{|\mathbf{G}|}\sum_{R \subset \mathbf{G}}\psi^{i}(R)\psi^{j}(R)\psi^{k}(R)^{*}.
\tag{7.4.3}
\end{equation*}
当 $\mathbf{G}$ 是晶体空间群时，如何由该公式得到 $c_{ij, k}$ 的详细分析已在 4.7 节用小群方法给出。

本节证明了：两个不可约共表示的直积可以约化为若干不可约共表示之和，其系数可以完全由幺正子群内部的 Clebsch--Gordan 系数表达；三种情形 (a)、(b)、(c) 的组合共六种配对方式，各自的约化公式在原书中逐一给出。
% 7.5 磁点群的共表示
\section{磁点群的共表示}

\noindent\textbf{灰群；时间反演不变性的后果}\ 已提到过，时间反演操作的效果是反转原子或离子的磁矩。因此对顺磁或抗磁晶体中的离子，在没有外磁场时，晶体在时间反演操作下不变；于是可以说，普通点群（即 I 型 Shubnikov 点群）不足以描述该离子局部环境的对称性，更好的描述是相应的 II 型 Shubnikov 群之一——灰群。考虑这些群并导出其共表示、进而考察晶体在时间反演操作下的不变性是否引入额外简并，是相当直接了当的（Cracknell 1966\textit{a}, Cracknell and Wong 1967, Dimmock and Wheeler 1962\textit{b}, 1964）。

对 32 个灰群中的每一个，可取 $A$ 为元素 $\theta$（时间反演操作本身）。于是 $\mathbf{G} = \mathbf{M}$、$\mathbf{G}$ 在 $\mathbf{M}$ 中指数为 2，情形属于 (a)、(b) 或 (c) 视 $\Gamma^{i*} \simeq \Gamma^{i}$ 与否：对实表示（(b) 类），时间反演把每个态映到同一能级的简并伙伴（但不产生新的额外简并，除非表示为复的）；对复表示（(c) 类），${}^{1}\Gamma$ 与 ${}^{2}\Gamma$ 这一共轭对在时间反演下互相简并。这就是著名的 Kramers 简并：对含奇数个电子的体系，时间反演不变性保证至少二重简并。

\noindent\textbf{III 型磁点群的共表示}\ 对 58 个黑白点群，逐一选取幺正子群 $\mathbf{G}$、陪集代表 $A$，按情形 (a)/(b)/(c) 由表 2.2 的单值与双值表示导出共表示，并标明时间反演引起的额外简并。原书第 636--650 页逐一列表；以下为扫描。

\begin{figure}[H]
\centering
\includegraphics[width=0.86\textwidth]{c7_p649.jpg}
\caption*{\textbf{7.5 节的共表示表}：磁点群的不可约共表示（原书第 636--650 页扫描）}
\end{figure}
\begin{figure}[H]
\centering
\includegraphics[width=0.86\textwidth]{c7_p650.jpg}
\caption*{原书第 637 页}
\end{figure}
\begin{figure}[H]
\centering
\includegraphics[width=0.86\textwidth]{c7_p651.jpg}
\caption*{原书第 638 页}
\end{figure}
\begin{figure}[H]
\centering
\includegraphics[width=0.86\textwidth]{c7_p652.jpg}
\caption*{原书第 639 页}
\end{figure}
\begin{figure}[H]
\centering
\includegraphics[width=0.86\textwidth]{c7_p653.jpg}
\caption*{原书第 640 页}
\end{figure}
\begin{figure}[H]
\centering
\includegraphics[width=0.86\textwidth]{c7_p654.jpg}
\caption*{原书第 641 页}
\end{figure}
\begin{figure}[H]
\centering
\includegraphics[width=0.86\textwidth]{c7_p655.jpg}
\caption*{原书第 642 页}
\end{figure}
\begin{figure}[H]
\centering
\includegraphics[width=0.86\textwidth]{c7_p656.jpg}
\caption*{原书第 643 页}
\end{figure}
\begin{figure}[H]
\centering
\includegraphics[width=0.86\textwidth]{c7_p657.jpg}
\caption*{原书第 644 页}
\end{figure}
\begin{figure}[H]
\centering
\includegraphics[width=0.86\textwidth]{c7_p658.jpg}
\caption*{原书第 645 页}
\end{figure}
\begin{figure}[H]
\centering
\includegraphics[width=0.86\textwidth]{c7_p659.jpg}
\caption*{原书第 646 页}
\end{figure}
\begin{figure}[H]
\centering
\includegraphics[width=0.86\textwidth]{c7_p660.jpg}
\caption*{原书第 647 页}
\end{figure}
\begin{figure}[H]
\centering
\includegraphics[width=0.86\textwidth]{c7_p661.jpg}
\caption*{原书第 648 页}
\end{figure}
\begin{figure}[H]
\centering
\includegraphics[width=0.86\textwidth]{c7_p662.jpg}
\caption*{原书第 649 页}
\end{figure}
\begin{figure}[H]
\centering
\includegraphics[width=0.86\textwidth]{c7_p663.jpg}
\caption*{原书第 650 页}
\end{figure}

% 7.6 磁空间群的共表示
\section{磁空间群的共表示}

确定磁空间群共表示的问题与确定磁点群共表示的问题之间的关系，非常类似于由点群表示确定 Fedorov 群（普通空间群）表示的问题之间的关系。由点群表示确定普通空间群表示的理论已在 3.7 与 3.8 节考虑过，并在第四章更严格地处理。记得利用第四章的诱导表示理论，空间群表示的确定问题可以简化：对表示域中的每个波矢 $\mathbf{k}$，只需推导小陪群 $\overline{\mathbf{G}}^{\mathbf{k}}$ 的、带有依赖于该 $\mathbf{k}$ 矢量的因子系统的射影表示。类似地，磁空间群的共表示可以通过诱导从小群的共表示得到；II 型（灰）空间群的共表示由相应 Fedorov 群表示的 (a)/(b)/(c) 分类直接给出，III 型（黑白）空间群则需对每个 $\mathbf{k}$ 辨认使其翻转为 $-\mathbf{k}$ 的反幺正元素，并把共表示按三种情形构造。原书第 652--664 页对磁空间群逐一列表；以下为扫描。

\begin{figure}[H]
\centering
\includegraphics[width=0.86\textwidth]{c7_p665.jpg}
\caption*{\textbf{7.6 节的共表示表}：磁空间群的不可约共表示（原书第 652--664 页扫描）}
\end{figure}
\begin{figure}[H]
\centering
\includegraphics[width=0.86\textwidth]{c7_p666.jpg}
\caption*{原书第 653 页}
\end{figure}
\begin{figure}[H]
\centering
\includegraphics[width=0.86\textwidth]{c7_p667.jpg}
\caption*{原书第 654 页}
\end{figure}
\begin{figure}[H]
\centering
\includegraphics[width=0.86\textwidth]{c7_p668.jpg}
\caption*{原书第 655 页}
\end{figure}
\begin{figure}[H]
\centering
\includegraphics[width=0.86\textwidth]{c7_p669.jpg}
\caption*{原书第 656 页}
\end{figure}
\begin{figure}[H]
\centering
\includegraphics[width=0.86\textwidth]{c7_p670.jpg}
\caption*{原书第 657 页}
\end{figure}
\begin{figure}[H]
\centering
\includegraphics[width=0.86\textwidth]{c7_p671.jpg}
\caption*{原书第 658 页}
\end{figure}
\begin{figure}[H]
\centering
\includegraphics[width=0.86\textwidth]{c7_p672.jpg}
\caption*{原书第 659 页}
\end{figure}
\begin{figure}[H]
\centering
\includegraphics[width=0.86\textwidth]{c7_p673.jpg}
\caption*{原书第 660 页}
\end{figure}
\begin{figure}[H]
\centering
\includegraphics[width=0.86\textwidth]{c7_p674.jpg}
\caption*{原书第 661 页}
\end{figure}
\begin{figure}[H]
\centering
\includegraphics[width=0.86\textwidth]{c7_p675.jpg}
\caption*{原书第 662 页}
\end{figure}
\begin{figure}[H]
\centering
\includegraphics[width=0.86\textwidth]{c7_p676.jpg}
\caption*{原书第 663 页}
\end{figure}
\begin{figure}[H]
\centering
\includegraphics[width=0.86\textwidth]{c7_p677.jpg}
\caption*{原书第 664 页}
\end{figure}

% 7.7 例子：P4'2/mnm' 的空间群共表示及其 Kronecker 积
\section{例子：$P4'_{2}/mnm'$ 的空间群共表示及其 Kronecker 积}

作为前几节理论应用的说明，我们考虑推导 III 型 Shubnikov 空间群 $P4'_{2}/mnm'$ 的共表示；这是反铁磁体 $\mathrm{MnF}_{2}$ 的空间群。该空间群曾由 Dimmock and Wheeler (1962\textit{b}) 与 Cracknell (1967\textit{c}) 考虑过。结构示于图 7.5，空间群的元素可由表 7.2 与 3.7 辨认。显然它关联于普通空间群 $P4_{2}/mnm\ (D^{14}_{4h})$；由表 7.2 可见其生成元与 $P4_{2}/mnm$ 相同，只是 $\{C^{+}_{4z} \mid {\frac{1}{2}}{\frac{1}{2}}{\frac{1}{2}}\}$ 必须乘以 $\theta$。处理 $\mathrm{MnF}_{2}$ 的结构时，用《国际表》给出的原点与轴向设定是方便的。原书第 666--689 页逐一推导了各对称点的共表示（按 (a)/(b)/(c) 情形）以及它们的 Kronecker 积约化；以下为扫描。

\begin{figure}[H]
\centering
\includegraphics[width=0.86\textwidth]{c7_p679.jpg}
\caption*{\textbf{7.7 节的表}：$P4'_{2}/mnm'$ 的共表示及其直积（原书第 666--689 页扫描）}
\end{figure}
\begin{figure}[H]
\centering
\includegraphics[width=0.86\textwidth]{c7_p680.jpg}
\caption*{原书第 667 页}
\end{figure}
\begin{figure}[H]
\centering
\includegraphics[width=0.86\textwidth]{c7_p681.jpg}
\caption*{原书第 668 页}
\end{figure}
\begin{figure}[H]
\centering
\includegraphics[width=0.86\textwidth]{c7_p682.jpg}
\caption*{原书第 669 页}
\end{figure}
\begin{figure}[H]
\centering
\includegraphics[width=0.86\textwidth]{c7_p683.jpg}
\caption*{原书第 670 页}
\end{figure}
\begin{figure}[H]
\centering
\includegraphics[width=0.86\textwidth]{c7_p684.jpg}
\caption*{原书第 671 页}
\end{figure}
\begin{figure}[H]
\centering
\includegraphics[width=0.86\textwidth]{c7_p685.jpg}
\caption*{原书第 672 页}
\end{figure}
\begin{figure}[H]
\centering
\includegraphics[width=0.86\textwidth]{c7_p686.jpg}
\caption*{原书第 673 页}
\end{figure}
\begin{figure}[H]
\centering
\includegraphics[width=0.86\textwidth]{c7_p687.jpg}
\caption*{原书第 674 页}
\end{figure}
\begin{figure}[H]
\centering
\includegraphics[width=0.86\textwidth]{c7_p688.jpg}
\caption*{原书第 675 页}
\end{figure}
\begin{figure}[H]
\centering
\includegraphics[width=0.86\textwidth]{c7_p689.jpg}
\caption*{原书第 676 页}
\end{figure}
\begin{figure}[H]
\centering
\includegraphics[width=0.86\textwidth]{c7_p690.jpg}
\caption*{原书第 677 页}
\end{figure}
\begin{figure}[H]
\centering
\includegraphics[width=0.86\textwidth]{c7_p691.jpg}
\caption*{原书第 678 页}
\end{figure}
\begin{figure}[H]
\centering
\includegraphics[width=0.86\textwidth]{c7_p692.jpg}
\caption*{原书第 679 页}
\end{figure}
\begin{figure}[H]
\centering
\includegraphics[width=0.86\textwidth]{c7_p693.jpg}
\caption*{原书第 680 页}
\end{figure}
\begin{figure}[H]
\centering
\includegraphics[width=0.86\textwidth]{c7_p694.jpg}
\caption*{原书第 681 页}
\end{figure}

% 7.8 自旋空间群
\section{自旋空间群}

本章描述的磁点群或磁空间群的概念，基于这样一个假设：算符同时作用于粒子的空间坐标与自旋坐标。这样的操作若要是晶体的对称操作，其结果是原子位置与自旋从初始值起不可分辨。然而，若对磁晶体的自旋哈密顿量的形式作某些简化假定，它可能具有比磁空间群所描述的更多的对称性。于是可以定义并使用\textit{自旋空间群}（spin-space group）$\mathbf{G}_{s}$，它比晶体的磁空间群具有相当多的对称性（Brinkman 1967, Brinkman and Elliott 1966\textit{a}, \textit{b}）。哈密顿量精确形式的细节、连同所涉近似之物理意义的考察，由 Brinkman and Elliott 给出。例如自旋波色散关系的总体特征随后便可确定。自旋空间群是空间群（作用于空间坐标）与自旋转动群（作用于自旋坐标）的半直积子群；其不可约表示可由两个因子的表示构造，且自旋波分支可按其标记。

% 7.9 多色点群与空间群
\section{多色点群与空间群}

Shubnikov 群（黑白群）是通过在位置坐标 $x, y, z$ 之外添加一个取两个可能值的额外坐标 $s$ 而导出的。若让额外坐标 $s$ 不再是二值而是三值、四值或六值的，这一做法可以进一步推广。正如二色点群与空间群的推导，可以推导三色、四色或六色点群与空间群。一般地考虑 $p$ 色点群，$p = 2, 3, 4$ 或 6。设我们有按某种顺序排列的 $p$ 种颜色，例如三色红、绿、黄，则可以定义换色操作 $\mathcal{R}_{p}$，它把这 $p$ 种颜色循环置换，例如 $\mathcal{R}_{3}$ 把红变为绿、绿变为黄、黄变为红（见图 7.10）。$p$ 色点群将有三种类型。I 型 $p$ 色点群（或空间群）就是普通的点群（或空间群）$\mathbf{G}$。II 型 $p$ 色群含有 $\mathcal{R}_{p}$ 的全部幂（至 $p$ 次），从而是灰群对 $p$ 色的推广；III 型 $p$ 色群则含 $\mathcal{R}_{p}$ 的部分幂与 $\mathbf{G}$ 的部分元素之积。$p$ 色群的数目可以由群论论证推出；对 $p = 3, 4, 6$ 的清单与讨论见原书本节（含文献）。多色群在磁结构（多子格磁性）与振动光谱等问题中已有应用。

